% from main.tex
%%%%%%%%%%  scan idx 386  |  folio 373  %%%%%%%%%%
aient des degrés compris entre \(n_0\) et \(n_1\). On procède alors par récurrence sur \(N=n_1-n_0\).

Lorsque \(N=0\), i.e. \(n_1=n_0\), soit \(n\) la valeur commune. L'homomorphisme canonique (1.2.1)
\[
  M_n\otimes_{\mathcal O_S}B(-n)\longrightarrow M
\]
est alors un isomorphisme. En effet, la vérification est locale, et d'après (1.2.3), on peut alors supposer que \(M\simeq(B(-n))^k\), et le résultat est trivial.

Dans le cas \(N>0\), on considère encore l'homomorphisme (1.2.1) :
\[
  M_{n_0}\otimes_{\mathcal O_S}B(-n_0)\longrightarrow M.
\]
On vérifie localement que cet homomorphisme est injectif, et que son conoyau \(M'\) est localement libre de type fini, à générateurs locaux de degré compris entre \(n_0+1\) et \(n_1\) : en effet, au voisinage de tout point de \(S\), \(M_{n_0}\) est soit nul, soit la composante homogène non nulle de \(M\) dont le degré est minimum; utilisant une décomposition de \(M\) en facteurs directs du type (1.2.2), on voit dans ce dernier cas en utilisant (1.2.3) que
\[
  M_{n_0}\otimes_{\mathcal O_S}B(-n_0)
\]
est isomorphe au facteur direct \((B(-n_0))^{k_{n_0}}\) de \(M\), d'où les propriétés annoncées. On applique alors l'hypothèse de récurrence à \(M'\), d'où aussitôt la conclusion voulue.

\textbf{1.8.} Revenant à la démonstration du théorème 1.1, montrons que \(K^{\bullet}(X)\) est engendré comme \(K^{\bullet}(S)\)-module par les puissances de \(\xi\).

Si \(F\) est un \(\mathcal O_X\)-Module localement libre de type fini, il existe d'après 1.4 un \(B\)-Module gradué de présentation finie \(M\), plat sur \(S\), tel que \(F=\Projz(M)\). D'après 1.6, \(M\) admet une résolution graduée par des \(B\)-Modules gradués localement libres de type fini. On en déduit une
